\documentclass[12pt]{article}

\usepackage {preamble}

\begin{document}

\author{}
\title{Системы управления базами данных}
\date{5 семестр, 2026}
\maketitle

Дейт <<Введение в СУБД>>

База данных "--- набор постоянных данных, которые используются прикладными системами для какого"=либо предприятия.

СУБД "--- программно"=аппаратный комплекс, обеспечивающий сохранность данных, целостность, доступ к данным.

Первичный ключ "--- уникальный идентификатор.

Домен "--- допустимое потенциальное множество значений данного типа.

Свойства отношений:

\begin{itemize}
	\item Все значения атрибутов атомарны (в отличие от таблиц).
	\item Кортежи и атрибуты не упорядочены.
	\item В отношении нет одинаковых кортежей.
\end{itemize}

\textbf{Первая нормальная форма.} "--- все атрибуты отношения содержат только атомарные значения.

Отношение $<=>$ 1НФ

Целостность: целостность сущностей и целостность внешних ключей

NULL"=значение "--- маркер, что значение неизвестно.

Парадокс трёхзначной логики. Неверно, что null = null; неверно, что null != null; a or (not a) не обязательно истина. 

Первичный ключ не имеет права быть null"=значением.

Подмножество млжет быть потенциальным ключом, если оно уникально и неизбыточно (никакое его подмножество не обладает свойством уникальности).

Первичным ключом могут быть числа или специальные типы данных, предназначенные для их генерации.

Связи: один ко многим, один к одному, многое ко многим (реализуется через несколько связей один ко многим)

set null и set default.

Стратегии: CASCADE (удаляем всю родительскую цепочку) И RESTRICT (если что"=то может пойти не так, запрещаем).

RESTRICT: не разрешать выполнение операции, приводящей к нарушению ссылочной целостности (от родителя к ребёнку).

CASCADE: разрешить выполнение требуемой операции, но изменить при этом другие отношения, чтобы не нарушить ссылочную целостность.

Автокаскад не стоит использовать. Лучше не удалять, а поставить статус <<неактивен>>.

Триггер "--- автоматически вызываемая функция.

Реляционная алгебра. Язык SQL "--- стандарт для доступа к реляционным данным. Реляционно полный язык: поддерживает алгебраические операции объединения, пересечения, вычитания и декартова произведения; специальные реляционные операции выборки (SELECT), проекции, соединения (JOIN: left join, right join, inner join), деления.

Предметная область "--- часть реального мира, данные из которой мы хотим отразить в базе данных.

Модель предметной области "--- наши знания о предметной области.

Физическая модель данных "--- привязка к конкретной СУБД.

\begin{enumerate}
\item Сама предметная область
\item Модель предметной области
\itrm Логическая модель данных
\item Физическая модель данных
\item Сама база данных и приложения
\end{enumerate}

Отношение находится во второй нормальной форме (2НФ), когда оно находится в 1НФ и нет неключевых атрибутов, зависящих от части сложного ключа. Неключевой атрибут "---  атрибут, не входящий в состав никакого потенциального ключа.

Атрибуты называются взаимно независимыми, если ни один из них функционально не зависим от другого.

Отношение находится в третьей нормальной форме (3НФ), когда оно находится в 2НФ и все неключевые атрибуты взаимно независимы.

Не стоит вешать индексы на время, строки, вещественные числа. Идеальные индексы "--- целые числа и перечислимые типы данных.

Селективность индекса. Если известно, что в этом наборе полей мало различных полей по сравнению с общим объёмом таблицы, нет смысла вешать туда индекс.

Фильтруемый индекс.

Транзакция "--- неделимая единица.

АСИТ: атомарность, согласованность, изоляция, долговечность.

Декларативная/процедурная поддержка ограничений целостности.

Протокол двухфазной блокировки:
\begin{enumerate}
	\item Перед выполнением каких"=либо операций с некоторым объектом транзакция должна заблокировать этот объект.
	\item После снятия блокировки транзакция не должна накладывать никаких других блокировок.
\end{enumerate}

\textbf{Теорема Есворана} "--- Если все транзакции в смеси подчиняются протоколу двухфазной блокировки, то для всех чередующихся графиков запуска существует возможность упорядочения.

Механизм репликации.

\begin{itemize}
	\item snapshot (репликация моментальных снимков);
	\item transactional (репликация транзакций);
	\item merge (репликация слиянием);
\end{itemize}
\end{document}

